\documentclass[11pt]{article}
\usepackage[margin=1in]{geometry}
\usepackage{booktabs}
\usepackage{hyperref}

\title{CAI6734 Weekly Report 5: Evaluating the Trained LA-CDM Model}
\author{Omer Kahveci}
\date{October 27, 2026}

\begin{document}
\maketitle

\section{Project Objective}

This week I focused on evaluating the Meta-Llama-3-8B-Instruct LA-CDM model
after its 300-step training cycle. The immediate goal is a matched comparison
between the trained model and the untrained Llama baseline on the same
240-case MIMIC-CDM test split, using one rollout per case.

\section{Progress}

I evaluated the saved LoRA adapter from the completed training run on the
fixed MIMIC-CDM test split. Both the trained and untrained Llama evaluations
used the same 240 cases and one rollout per case.

The untrained baseline produced a valid final diagnosis for 29 of 240 cases.
Its overall diagnosis accuracy was 8.75\%, counting undiagnosed cases as
incorrect, while accuracy among completed diagnoses was 72.41\%. Its
undiagnosed fraction was 87.92\%. This gap suggests that action-protocol
completion is a major limitation of the untrained model.

The completed training run also produced the validation results in
Table~\ref{tab:validation}. This pass used 240 validation cases with four
rollouts per case, for 960 predictions. It is reported separately from the
matched test-set comparison in Table~\ref{tab:test}.

\begin{table}[ht]
  \centering
  \caption{Trained Llama validation results after 300 training steps.}
  \label{tab:validation}
  \begin{tabular}{lr}
    \toprule
    Metric & Trained validation \\
    \midrule
    Diagnosis accuracy, all predictions & 56.46\% \\
    Diagnosis accuracy, diagnosed only & 67.50\% \\
    Undiagnosed fraction & 16.35\% \\
    Action-format reward & 0.8365 \\
    Macro F1, all predictions & 57.59\% \\
    Hypothesis accuracy & 63.06\% \\
    Expected calibration error & 7.89\% \\
    \bottomrule
  \end{tabular}
\end{table}

\begin{table}[ht]
  \centering
  \caption{Llama results on the same 240-case test split, with one rollout per
  case.}
  \label{tab:test}
  \begin{tabular}{lrr}
    \toprule
    Metric & Untrained test & Trained test \\
    \midrule
    Diagnosis accuracy, all predictions & 8.75\% & 68.75\% \\
    Diagnosis accuracy, diagnosed only & 72.41\% & 70.21\% \\
    Undiagnosed fraction & 87.92\% & 2.08\% \\
    Action-format reward & 0.1208 & 0.9792 \\
    Macro F1, all predictions & 15.54\% & 66.27\% \\
    Hypothesis accuracy & 65.25\% & 63.26\% \\
    Expected calibration error & 10.87\% & 10.38\% \\
    Average tests requested & 1.16 & 2.73 \\
    Average diagnostic cost & \$798.08 & \$1,137.67 \\
    \bottomrule
  \end{tabular}
\end{table}

\section{Published LA-CDM Reference}

Table~\ref{tab:paper} records the trained Qwen-2.5-7B-Instruct results
reported by Bani-Harouni et al. in Table~1 and Section~6 of
\href{https://arxiv.org/abs/2506.13474}{the LA-CDM paper}. These results
provide a reference for the reproduction, but differences in the base model,
training duration, and evaluation setup prevent a controlled comparison with
the Llama results above. In particular, the paper's mean class accuracy is
not the same metric as our overall accuracy.

\begin{table}[ht]
  \centering
  \caption{Published trained Qwen LA-CDM test results.}
  \label{tab:paper}
  \begin{tabular}{lr}
    \toprule
    Metric & Paper result \\
    \midrule
    Mean class accuracy & 81.3\% \\
    Micro F1 & 84.1\% \\
    Macro F1 & 81.3\% \\
    Average test cost & \$1,295.61 \\
    Hypothesis accuracy & 81.9\% \\
    Expected calibration error & 3.7\% \\
    \bottomrule
  \end{tabular}
\end{table}

\section{Interpretation and Next Steps}

The matched test results show a 60.00-percentage-point increase in overall
diagnosis accuracy: 21 of 240 cases were correct for the untrained model,
versus 165 of 240 for the trained model. The trained model completed a valid
diagnosis in 235 cases, leaving five undiagnosed. Accuracy among diagnosed
cases was 70.21\%, slightly below the untrained model's 72.41\% on its much
smaller set of 29 completed diagnoses. The main measured gain is therefore
reliable completion of the action protocol, rather than higher accuracy
conditional on completing a diagnosis.

The gain came with more test use: the average rose from 1.16 to 2.73 tests,
and estimated diagnostic cost rose from \$798.08 to \$1,137.67 per case.
Hypothesis accuracy decreased slightly from 65.25\% to 63.26\%, while expected
calibration error changed from 10.87\% to 10.38\%. The trained test evaluation
took 5,864.7 seconds, or about 98 minutes. The validation results above are
not directly comparable to the untrained test results because they used a
different split and four rollouts per case.

Next I will inspect representative incorrect and undiagnosed trajectories,
including the five cases without a final diagnosis, and examine whether
additional tests improved decisions enough to justify their cost.

\section{AI Disclaimer}

I used OpenAI Codex to help inspect the evaluation workflow and draft this
report. I reviewed the reported metrics and remain responsible for the final
interpretation.

\end{document}
